\section{Legal Properties}
\label{sec:legal}

\subsection{Procedural Fairness of Local ReCom Sampling}
\label{sec:legal:procedural}

BisectionEnsemble inherits the procedural fairness argument of ReCom-based
ensemble methods \citep{deford2021recombination, duchin2021report} while
preserving the structural guarantees of the bisection tree
\citep{dellaLibera2026apportion}. We address each legal property in turn.

\paragraph{No partisan inputs.}
The local ReCom chain at each bisection node operates on the census-tract
adjacency graph with TIGER boundary weights. No partisan data enters the
chain: the spanning tree distribution is determined by graph structure
(adjacency and population) alone, and the selection criterion (edge-cut rank)
is a geometric property. This is identical to the non-partisan posture of
ApportionRegions \citep{dellaLibera2026apportion} and satisfies the
``political blindness'' standard for neutral map-drawing algorithms articulated
in state redistricting jurisprudence. Post-\emph{Rucho v.\ Common Cause}
\citep{rucho2019} (which held federal courts lack jurisdiction over partisan
gerrymandering claims), partisan gerrymandering claims are litigated in state
courts under state constitutional provisions; the operative authorities for
algorithmic neutrality are \emph{League of Women Voters v.\ Commonwealth}
(Pa.\ 2018), \emph{Common Cause v.\ Lewis} (N.C.\ 2019), and \emph{Harper v.\
Hall} (N.C.\ 2022), all of which articulate procedural neutrality standards
that BisectionEnsemble satisfies by virtue of using no partisan inputs.

\paragraph{Uniform local sampling.}
The uniform spanning tree (UST) used in each ReCom step provides a
theoretically grounded sampling distribution: the UST of a graph $H$ is the
unique maximum-entropy distribution over spanning trees of $H$
\citep{aldous1990random, wilson1996generating}. Each accepted plan
$(A_t, B_t)$ is drawn from the distribution induced by the UST on the
feasible bisections of the local subregion. This provides the same
``representativeness'' argument as full-state ReCom ensembles, localised to
each bisection node.

\paragraph{Pre-committed procedure.}
The hyperparameters $p$ and $T$ are fixed before any plans are generated.
The percentile $p$ is the same legal posture argument as in PercentileSweep
\citep{dellaLibera2026percentilesweep}: a value of $p = 0.5$ selects the
median plan (most representative of the local feasible space), while $p = 0.0$
selects the most compact plan. Either choice can be defended as a pre-committed,
politically neutral criterion. No post-hoc adjustment based on partisan outcomes
is made.

\subsection{Bisection Tree Structural Guarantees}
\label{sec:legal:structure}

The bisection tree structure provides two constitutional guarantees that
BisectionEnsemble inherits by construction.

\paragraph{Population balance.}
Each bisection node in the tree receives a population target $(t_1, t_2)$
derived from its parent's allocation and the seat count $k$. The ReCom chain
only accepts plans that achieve population balance within the tolerance
$\varepsilon$ (by Theorem~\ref{thm:tractable}, such plans always exist).
The hierarchical structure ensures that population imbalances do not compound:
an imbalance at node $v$ is bounded by $\varepsilon$ and propagates only to
$v$'s subtree, where sub-tolerances absorb it. The final districts satisfy
the \emph{Wesberry v.\ Sanders} standard \citep{wesberry1964} of $\pm 0.5\%$
population deviation from the ideal, which is enforced at the leaf level.

\paragraph{Contiguity.}
By Proposition~\ref{prop:contiguity}, every plan produced at every node is
contiguous. Since contiguity is preserved at every level of the recursion,
the final districts are contiguous. This satisfies the contiguity requirement
that courts have interpreted as implicit in Article~I \S2 and explicit in
most state redistricting statutes.

\subsection{Percentile Choice and Equal Protection}
\label{sec:legal:percentile}

The choice of percentile $p$ has been challenged in redistricting litigation
as a ``hidden thumb on the scale.'' We address this concern directly.

\paragraph{$p = 0.5$ (median): the ``typical'' feasible plan.}
Selecting the median-edge-cut bisection at each node produces the plan that
is most representative of the local feasible bisection space. To the extent
the local chain mixes over the feasible bisection space of each node --- a
property whose formal proof is left to future work --- the distribution of
accepted plans approximates the stationary distribution over feasible bisections.
The median plan is thus an approximation to the plan that minimises the maximum
deviation from any other feasible bisection in the edge-cut metric. This is a
principled ``central tendency'' criterion that courts have accepted in the context
of full-state ensembles \citep{herschlag2020quantifying}. Empirically, the
acceptance rates observed (NC: 61\%, WI: 84\%, TX at bisection nodes: $>50\%$)
suggest reasonable mixing at $T=100$ for these states.

\paragraph{$p = 0.0$ (minimum): compactness maximisation.}
Selecting the minimum-edge-cut bisection at each node is equivalent to
the ReCom-constrained version of minimum-edge-cut redistricting
\citep{dellaLibera2026mec}: the most compact feasible bisection of each
subregion. The legal argument is that compactness is an explicitly enumerated
redistricting criterion in 37 states and has been accepted as a non-partisan
tie-breaker in judicial review of redistricting plans.

\paragraph{Sensitivity to $p$.}
Rather than arguing retrospectively that NC's stable result proves neutrality,
we commit to the following prospective sensitivity-analysis protocol: for any
state where BisectionEnsemble is deployed in litigation, run the full sweep
$p \in \{0.0, 0.25, 0.5, 0.75, 1.0\}$ and disclose the distribution of seat
counts across $p$ values. A result that is stable across $p$ is evidence of
neutrality; a result that varies with $p$ must be disclosed as such. This
commitment is more legally robust than a retrospective stability claim, because
it applies regardless of whether the target state resembles NC. As shown in
Table~\ref{tab:partisan}, WI exhibits a one-seat shift at $p=0.5$ relative
to $p=0.0$; under this protocol, this variation would be disclosed and its
geographic cause (the Rodden sorting effect, see Section~\ref{sec:results:partisan})
identified, rather than suppressed.

\paragraph{METIS seed protocol.}
The METIS seed for each node is drawn from the system CSPRNG (cryptographically
secure pseudorandom number generator) before any plans are computed. The seed is
logged before the first plan is generated at each node. If the METIS seed were
chosen after inspecting results, the ``pre-committed'' argument would break down;
the CSPRNG-draw-before-compute protocol ensures the seed is not a post-hoc choice.

\paragraph{Disclosure requirements.}
Any deployment of BisectionEnsemble in litigation or administrative redistricting
should disclose (1) the value of $p$, (2) the value of $T$, (3) the METIS seed
used for initialisation (drawn via CSPRNG before any plans are generated),
(4) the census release identifier, (5) the acceptance rate at each node, and
(6) the full $p$-sweep seat distribution for the target state. Optionally, the
full sequence of accepted plans $\Pi$ (or their edge-cut values and population
assignments) may be preserved with the \texttt{--log-ensemble} flag, enabling
independent verification of the percentile selection. All items are logged by
the \texttt{redist-cli} implementation in the canonical JSONL audit log,
satisfying the deposition disclosure requirements of \citet{dellaLibera2026depo}.
